\documentclass{article}

\usepackage[a4paper]{geometry}
\usepackage[english]{babel}
\usepackage{parskip}

\usepackage{amsmath, amssymb}
\usepackage{graphicx}
\usepackage{subcaption}
\usepackage{float}
\usepackage[hidelinks]{hyperref}

\usepackage{listings}
\usepackage{xcolor}

\definecolor{codegreen}{rgb}{0,0.6,0}
\definecolor{codegray}{rgb}{0.5,0.5,0.5}
\definecolor{codepurple}{rgb}{0.58,0,0.82}
\definecolor{backcolour}{rgb}{0.95,0.95,0.92}

\lstdefinestyle{mystyle}{
  backgroundcolor=\color{backcolour},
  commentstyle=\color{codegreen},
  keywordstyle=\color{magenta},
  numberstyle=\tiny\color{codegray},
  stringstyle=\color{codepurple},
  basicstyle=\ttfamily\footnotesize,
  breaklines=true,
  captionpos=b,
  keepspaces=true,
  numbers=left,
  numbersep=5pt,
  showstringspaces=false,
  tabsize=2
}
\lstset{style=mystyle}

\newcommand{\CodeDir}{Code}
\newcommand{\PlotDir}{Plots}
\newcommand{\CalcDir}{Calculations}

\title{NURa Hand-in 3}
\author{Name (Student number)}
\date{Date}

\begin{document}
\maketitle

\textit{Write a short description of your solution in each section. Explain your choices. Explicitly draw conclusions from your results.}

\textit{You can refer to selected lines of code where you think it is necessary. However, you must provide the full code relevant to each (sub)question below it.}

\textit{Don't forget to describe your figures and mention what conclusions you draw from them!}


\section{Satellite galaxies around a massive central}

\subsection{(1a) Maximum of N(x)}

Our goal is to find the radius $x\equiv r/r_\mathrm{vir}$ at which we expect to find the most satellite galaxies.
This comes down to finding the maximum of
\begin{equation}
N(x) = 4\pi x^2 n(x),
\end{equation}
where $n(x)$ is the number density of satellite galaxies given by
\begin{equation}
n(x) = A \langle N_\mathrm{sat} \rangle 
\left( \frac{x}{b}\right)^{a-3}
\exp\left[-\left(\frac{x}{b}\right)^c\right].
\end{equation}
First we inspect the function of interest by plotting it over the range $x \in [10^-4, 5]$ (see Fig. \ref{Fig: Satellite_number_plot}).
\begin{figure}[H]
\centering
\includegraphics[width=0.85\textwidth]{\PlotDir/1a_N(x).png}
\caption{Expected number of satellites at a given radius x.}
\label{Fig: Satellite_number_plot}
\end{figure}
We see that it has a maximum somewhere in between $x=0.1$ and $x=1$, while not containing any other local maxima. 
In order to find the maximum, we employ a golden section search that \textbf{minimizes} the negative of the function of interest. 
Golden section search iteratively tightens a bracket in a self-similar fashion. 
A bracket is a set of 3 x-values where the middle point attains a lower value than the two edge points.
This algorithm also requires a bracket as input. 
Although we could construct our own bracket by inspecting Fig. \ref{Fig: Satellite_number_plot} closer, 
we instead use a golden ratio bracketing algorithm implemented in the Minimizer class (see code snippet).
This also speeds up the code in general, as it starts off with a bracket where 
the ratio of distances between the 3 x-values is the golden ratio.

The maximum of \(N(x)\) is found at:
$$x_{\text{max}} \approx  \input{\CalcDir/satellite_max_x.txt} $$
$$N(x_{\text{max}}) \approx  \input{\CalcDir/satellite_max_Nx.txt} $$

\subsubsection{Code for minimizer class}
\lstinputlisting[language=Python]{\CodeDir/minimizer.txt}

\subsubsection{Code for N(x) maximization}
NOTE: all code is reproduced below as an example, show only the relevant lines for this exercise here instead.
\lstinputlisting[language=Python]{\CodeDir/satellites_maximize_code.txt}

\subsection{(1b) $\chi^2$-fit}

Argue the choices of your binning. Explain the minimization routine used. Bonus points for a correctly implemented derivative-based routine.

\begin{table}[H]
\centering
\caption{Best-fit parameters for $\chi^2$ fits.}
\begin{tabular}{|c|c|c|c|c|c|}
\hline
Dataset & $N_{\mathrm{sat}}$ & $\chi^2$ & $a$ & $b$ & $c$ \\
\hline
\input{Calculations/table_fitparams_chi2.tex} \\
\hline
\end{tabular}
\end{table}

\begin{figure}[H]
\centering
\includegraphics[width=0.85\textwidth]{\PlotDir/satellite_fits_chi2.png}
\caption{Fit of the satellite number density profile using the chi-squared method.}
\end{figure}

\subsubsection{Code for $\chi^2$-fit}
NOTE: all code is reproduced below as an example, show only the relevant lines for this exercise here instead.
\lstinputlisting[language=Python]{\CodeDir/satellites_chi2_code.txt}
 
\subsection{(1c) Poisson ln-likelihood fit}

Explain how you determine the Poisson ln-likelihood. Describe the optimization method used. 
How do your best fit parameters differ from the ones obtained in (1b)? Why?

We want to find parameters $p = (a, b, c)$, such that the likelhihood $\mathcal{L}(p)$ is maximized.
We then bin the data into $N$ bins centered on radii $x_i$, and label the number of 
satellites in bin $x_i$ as $y_i$. The number of sattelites in each bin are then 
distributed according to a poisson distribution around some expected number predicted 
by a true model parametrized by parameters $p$; $\mu(x_i|p)$. 
So the probability of finding a $y_i$ satellites in bin $x_i$,
given a model $\mu(x_i|p)$ is given by
\begin{equation}
  \frac{
    \mu(x_i|p)^{y_i}e^{-\mu(x_i|p)}
  }
  {
    y_i!
  }
  .
\end{equation}

We assume that the datapoints of radii are sampled independently. 
Maximizing the likelihood is equivalent to minimizing the negative log-likehood which is given by:
\begin{equation}
  -\ln \mathcal{L}(p)
  =
  -\sum_{i=0}^{N-1} 
  \left[
    y_i \ln (\mu(x_i|p)) 
    - \frac{\mu(x_i|p)}
    - \ln(y_i!)
  \right].
\end{equation}
Writing out the sum this becomes
\begin{equation}
  -\ln \mathcal{L}(p)
  =
  -\sum_{i=0}^{N-1} 
    y_i \ln (\mu(x_i|p)) 
  +
  \sum_{i=0}^{N-1} 
    \frac{\mu(x_i|p)}
  +
  \sum_{i=0}^{N-1} 
    \ln(y_i!).
\end{equation}
Since our goal is to minimize this function for certain $p$, we can ignore the
last term since it does not depend on $p$. But also note, that the second term 
is simply the expected total number of satellites in the data. So if we 
constrain our model $\mu(x_i|p)$, such that it always satisfies $\sum{i=0}^{N-1} = N_\mathrm{sat}$.
Then this becomes independent of $p$ as well.
Then we are left with minimizing the function $g(p)$ given by
\begin{equation}
  g(p) \equiv -\sum_{i=0}^{N-1} y_i \ln (\mu(x_i|p)).
\end{equation}

To avoid any bias in our choice of binning we can take the limit to an infinite number of bins.
To be more precise, we choose the number of bins so large, such that $y_i \in \{0, 1\}$
for all bins $x_i$.
$g(p)$ can then be rewritten to:
\begin{equation}
  g(p)
  =
  -\sum_{i} \ln (\mu(x_i|p)).
\end{equation}
Where $i$ now sums over all the data points and $x_i$ is the value of each data point.


\begin{table}[H]
\centering
\caption{Best-fit parameters for $\chi^2$ fits.}
\begin{tabular}{|c|c|c|c|c|}
\hline
Dataset  & llh & $a$ & $b$ & $c$ \\
\hline
\input{Calculations/table_fitparams_poisson.tex} \\
\hline
\end{tabular}
\end{table}

\begin{figure}[H]
\centering
\includegraphics[width=0.8\textwidth]{\PlotDir/satellite_fits_poisson.png}
\caption{Fit of the satellite number density profile using the Poisson log-likelihood method.}
\end{figure}

\subsubsection{Code for Poisson ln-likelihood fit}
NOTE: all code is reproduced below as an example, show only the relevant lines for this exercise here instead.
\lstinputlisting[language=Python]{\CodeDir/satellites_poisson_code.txt}

\subsection{(1d) Statistical tests}

% change table to include KS score if you implemented it
\begin{table}[H]
\centering
\caption{Statistical tests for the $\chi^2$ and Poisson ln-likelihood fits.}
\begin{tabular}{|c|c|c|c|}
\hline
Method & Dataset & G & Q\\
\hline
\input{\CalcDir/statistical_test_table_rows.tex} \\
\hline
\end{tabular}
\end{table}

What can you conclude from the results of the statistical tests?

\subsubsection{Code for statistical tests}
NOTE: all code is reproduced below as an example, show only the relevant lines for this exercise here instead.
\lstinputlisting[language=Python]{\CodeDir/satellites_statistical_tests_code.txt}

\subsection{(1e) \textbf{Bonus} Monte Carlo simulation}

Pick a single dataset. Describe the implementation of your Monte Carlo parametric bootstrap. 

\begin{figure}[H]
\centering
\includegraphics[width=0.85\textwidth]{\PlotDir/satellite_monte_carlo_chi2.png}
\caption{Best $\chi^2$-fit profiles using pseudo-experiments. }
\end{figure}

\begin{figure}[H]
\centering
\includegraphics[width=0.85\textwidth]{\PlotDir/satellite_monte_carlo_poisson.png}
\caption{Best Poisson-likelihood fit profiles using pseudo-experiments. }
\end{figure}

\subsubsection{Code for Monte Carlo simulation}
NOTE: all code is reproduced below as an example, show only the relevant lines for this exercise here instead.
\lstinputlisting[language=Python]{\CodeDir/satellites_monte_carlo_code.txt}

\end{document}
